% !PRINTED VI
\weberpage{VI}
\noindent Although algebra in part goes beyond the theory of numbers, and reaches over into other fields, e.g.\ the theory of functions, or, in its applications, into geometry as well, the science of numbers nevertheless always remains the foremost example for all algebraic considerations, and the questions of number theory that stand today in the foreground of interest are predominantly of an algebraic nature. This marked out the path that I had to follow in my work.

The extensive material has been divided into two volumes. The first volume contains the elementary part of algebra, which, by a traditional expression, may be called \gesperrt{calculation with letters}, then the rules for the numerical computation of the roots of equations, and the beginnings of \gesperrt{Galois}'s theory.

The second volume, which will, I hope, follow the first within a short time, is to present the general theory of finite groups, the theory of linear substitution groups, and applications to various individual problems, and is to conclude with the theory of algebraic numbers, where an attempt is made to unite the various points of view from which this theory has hitherto been considered.

Finally, the second volume is to provide an alphabetical index to both volumes.

As cannot be otherwise in a discipline that is in the course of rapid development and on which work is being done from the most diverse quarters, in algebra, too, the use of language and the notation are very varied and frequently not in agreement. This makes a unified presentation, as well as penetration into the various works, very difficult.

I have therefore endeavoured to adhere uniformly to a mode of expression that is as appropriate as possible, and in doing so I have frequently taken counsel with colleagues. I may express the hope of having thereby contributed to the consolidation of a uniform terminology.

The references to the literature and the historical notes given in the book make no claim whatsoever to completeness, although I have endeavoured
